\documentclass[exercice]{thedocument}%
\usepackage{tkz-euclide}%
\startdatakeys
\source{}
\BO{2026}
\cycle{4}
\competencesDuSocle{RE}
\objectifsApprentissage{C4A3ee,C4B2ea,C4B1ee}
\automatismes{C4B3cc,C4B2db}
\enddatakeys
\begin{document}%
A, B et C sont trois points non alignés.
\begin{center}
    \begin{tikzpicture}
        \coordinate[label=below right:A](A)at(0,0);
        \coordinate[label=above right:B](B)at(5,0);
        \coordinate[label=left:C](C)at(60:5);
        \tkzDrawLines[add=0.2 and 0.2](A,C A,B C,B)
    \end{tikzpicture}
\end{center}
\begin{enumerate}
    \item Tracer la droite (d$_1$) perpendiculaire à (AB) passant par C.
    \item Tracer la droite (d$_2$) perpendiculaire à (BC) passant par A.
    \item Tracer la droite (d$_3$) perpendiculaire à (AC) passant par B.
    \item Que peut on dire de ces trois droites~?
\end{enumerate}
\end{document}
